\section{Prompts and Schemas}
\label{app:prompts}

All strategies share a single per-step policy prompt (\S\ref{app:rank}); the
reasoning-trajectory strategies add two prompts of their own
(\S\ref{app:reason}); both reward sources share one scoring prompt
(\S\ref{app:reward}). All calls request \texttt{response\_format
= json\_object} and are parsed defensively: malformed entries are dropped, and
any candidate missing from the output receives a floor score of $10^{-3}$ so
search never crashes on a parse failure. Code longer than 5{,}000 characters is
truncated with an explicit \texttt{/* \ldots truncated \ldots */} marker.

\subsection{Per-step ranking prompt (all taxonomy strategies)}
\label{app:rank}

Greedy, self-consistency, best-of-$N$, beam, and $\mctstax$ all call one shared
step: rank the candidate children of the current taxonomy node. Placeholders are
in angle brackets.

\begin{lstlisting}
SYSTEM:
You are a security expert classifying the vulnerability in a code
function according to the CWE taxonomy. Classify the WEAKNESS the
code exhibits, not the fix. You are given the candidate CWE
categories at one level of the taxonomy; choose among THESE
candidates only.

USER:
[Prior reasoning: <rationale carried over from the parent step>]

Function:
```c
<code, truncated at 5000 chars>
```

Candidate CWE categories:
<one line per child: "- CWE-XXX: <name>. <description>">

Score EVERY candidate listed above by how likely it describes the
weakness in the function (0.0 = no, 1.0 = yes); scores need not
sum to 1. You MUST output an entry for every candidate CWE id.
Return JSON: {"scores": {"CWE-XXX": 0.0, "CWE-YYY": 0.0, ...},
"rationale": "one sentence"}.
\end{lstlisting}

Greedy takes the arg-max at temperature $0$; self-consistency and best-of-$N$
sample the same prompt at temperature $0.7$ with a per-sample cache nonce;
beam and $\mctstax$ use the full score vector as child priors. The
\texttt{rationale} is capped at 300 characters and passed down as ``prior
reasoning'' to the next level.

\subsection{Reasoning-trajectory prompts ($\mctsreas$)}
\label{app:reason}

$\mctsreas$ searches over reasoning states rather than taxonomy nodes. A move is
either a brief next reasoning step or a commit to one leaf CWE; the candidate
list shown to the model is the flat list of tree leaves (id $+$ name).

\begin{lstlisting}
SYSTEM:
You are a security expert reasoning step by step about the weakness
in a code patch, in order to classify it into exactly one CWE.
Reason about the removed (vulnerable) lines, not the fix.

USER (expansion; k = 3):
Code patch:
```c
<code, truncated at 5000 chars>
```

Reasoning so far:
<numbered steps, or "(none yet)">

Candidate CWEs:
<one line per leaf: "- CWE-XXX: <name>">

Propose 3 DISTINCT next moves toward the answer. Each move is
either a brief next reasoning step, or a final commit to ONE
candidate CWE id. Return JSON: {"moves": [{"type": "step",
"text": "..."} or {"type": "commit", "cwe": "CWE-XXX"}, ...]}.

USER (forced commit; rollouts and final extraction):
<same code + reasoning-so-far + candidate blocks>

Commit to the single best CWE now. Return JSON {"cwe": "CWE-XXX"}.
\end{lstlisting}

Reasoning steps are capped at 300 characters and traces at
\texttt{max\_steps} $= 3$; a commit to an id outside the candidate list falls
back to the first leaf (and is counted as an error in practice). Expansion runs
at temperature $0$; rollout commits run at temperature $0.7$ with a
\texttt{(seed, iteration)} cache nonce.

\subsection{Reward prompt (self-evaluation and external judge)}
\label{app:reward}

Both reward sources use the \emph{same} prompt; the only difference is the model
it is sent to --- the policy itself (\texttt{self\_eval}) or the stronger judge
(\texttt{external\_judge}, Claude-Haiku-4.5). The reward scores a partial or
complete root-to-node path in $[0,1]$; both are zero-shot, no PRM is trained.

\begin{lstlisting}
SYSTEM:
You are auditing a proposed CWE classification of a code function.
Judge how well the proposed CWE path describes the weakness in the
code. Classify the weakness, not the fix.

USER:
Function:
```c
<code, truncated at 5000 chars>
```

Proposed CWE path: CWE-AAA (<name>) -> CWE-BBB (<name>) -> ...
This path is {PARTIAL (not yet a leaf) | COMPLETE (leaf)}.

Description of the final node CWE-BBB: <description or "n/a">

Return JSON {"score": 0.0-1.0} where score is the probability this
path correctly classifies the weakness in the function.
\end{lstlisting}

The score is clamped to $[0,1]$; a parse failure scores $0$.
